\documentclass[10pt,a4paper]{amsart}
\usepackage{amsmath,amssymb,amsthm,graphicx}
% Retained minimal preamble for the Stacks Project chapter more-algebra.tex (More on Algebra).
% The chapter file pinned-original-source.tex carries only \input{preamble} on its first line, so no
% preamble text lives inside the chapter itself.  The Stacks Project uses one single shared preamble for
% every chapter; the snapshot stages it identically as sources/stacks-chow/preamble.tex and
% sources/stacks-etale-cohomology/preamble.tex.  The block below is copied mechanically, line for line,
% from sources/stacks-chow/preamble.tex lines 14-18 (the Xy-pic setup) and 154-217 (the native theorem
% declarations together with the mathematical macros).  Nothing was added, removed or redefined; the
% wrapper class and the amsmath/amssymb/amsthm/graphicx base packages are supplied by the assembler.
% Chapter: more-algebra.tex -- More on Algebra.
% For commutative diagrams we use Xy-pic
\usepackage[all]{xy}

% We use 2cell for 2-commutative diagrams.
\xyoption{2cell}

% Theorem environments.
%
\theoremstyle{plain}
\newtheorem{theorem}[subsection]{Theorem}
\newtheorem{proposition}[subsection]{Proposition}
\newtheorem{lemma}[subsection]{Lemma}

\theoremstyle{definition}
\newtheorem{definition}[subsection]{Definition}
\newtheorem{example}[subsection]{Example}
\newtheorem{exercise}[subsection]{Exercise}
\newtheorem{situation}[subsection]{Situation}

\theoremstyle{remark}
\newtheorem{remark}[subsection]{Remark}
\newtheorem{remarks}[subsection]{Remarks}

\numberwithin{equation}{subsection}

% Macros
%
\def\lim{\mathop{\mathrm{lim}}\nolimits}
\def\colim{\mathop{\mathrm{colim}}\nolimits}
\def\Spec{\mathop{\mathrm{Spec}}}
\def\Hom{\mathop{\mathrm{Hom}}\nolimits}
\def\Ext{\mathop{\mathrm{Ext}}\nolimits}
\def\SheafHom{\mathop{\mathcal{H}\!\mathit{om}}\nolimits}
\def\SheafExt{\mathop{\mathcal{E}\!\mathit{xt}}\nolimits}
\def\Sch{\mathit{Sch}}
\def\Mor{\mathop{\mathrm{Mor}}\nolimits}
\def\Ob{\mathop{\mathrm{Ob}}\nolimits}
\def\Sh{\mathop{\mathit{Sh}}\nolimits}
\def\NL{\mathop{N\!L}\nolimits}
\def\CH{\mathop{\mathrm{CH}}\nolimits}
\def\proetale{{pro\text{-}\acute{e}tale}}
\def\etale{{\acute{e}tale}}
\def\QCoh{\mathit{QCoh}}
\def\Ker{\mathop{\mathrm{Ker}}}
\def\Im{\mathop{\mathrm{Im}}}
\def\Coker{\mathop{\mathrm{Coker}}}
\def\Coim{\mathop{\mathrm{Coim}}}

% Boxtimes
%
\DeclareMathSymbol{\boxtimes}{\mathbin}{AMSa}{"02}

%
% Macros for moduli stacks/spaces
%
\def\QCohstack{\mathcal{QC}\!\mathit{oh}}
\def\Cohstack{\mathcal{C}\!\mathit{oh}}
\def\Spacesstack{\mathcal{S}\!\mathit{paces}}
\def\Quotfunctor{\mathrm{Quot}}
\def\Hilbfunctor{\mathrm{Hilb}}
\def\Curvesstack{\mathcal{C}\!\mathit{urves}}
\def\Polarizedstack{\mathcal{P}\!\mathit{olarized}}
\def\Complexesstack{\mathcal{C}\!\mathit{omplexes}}
% \Pic is the operator that assigns to X its picard group, usage \Pic(X)
% \Picardstack_{X/B} denotes the Picard stack of X over B
% \Picardfunctor_{X/B} denotes the Picard functor of X over B
\def\Pic{\mathop{\mathrm{Pic}}\nolimits}
\def\Picardstack{\mathcal{P}\!\mathit{ic}}
\def\Picardfunctor{\mathrm{Pic}}
\def\Deformationcategory{\mathcal{D}\!\mathit{ef}}
\usepackage{cleveref}
\crefname{originallemma}{original lemma}{original lemmas}
\Crefname{originallemma}{Original lemma}{Original lemmas}
\makeatletter
\newlabel{lemma-compare-representatives-perfect}{{1}{source}{}{sourcepointer-1}{}}
\newlabel{lemma-compare-representatives-perfect@cref}{{[originallemma][0][]1}{[1][0][]source}{}{}{}}
\makeatother
\crefname{originallemma}{original lemma}{original lemmas}
\Crefname{originallemma}{Original lemma}{Original lemmas}
\makeatletter
\newlabel{lemma-ideals-generated-by-minors}{{2}{source}{}{sourcepointer-2}{}}
\newlabel{lemma-ideals-generated-by-minors@cref}{{[originallemma][0][]2}{[1][0][]source}{}{}{}}
\makeatother
\title{The Fitting-type ideal attached to a bounded complex of finite free modules is well defined up to powers of the nonzerodivisor depending on the Euler characteristic}
\author{The Stacks Project Authors}
\date{stacks-project:more-algebra.tex:local-snapshot-20261012; GNU FDL 1.2 or later}
\begin{document}
\maketitle
\setcounter{subsection}{1}
\noindent\textit{Editorial note.} All delivered mathematical text is the author's own, copied from the pinned Stacks Project chapter file more-algebra.tex; no statement or proof step was rewritten, reordered or removed.\par\medskip
\noindent\textit{Editorial note.} The chapter file pinned-original-source.tex carries only \textbackslash{}input\{preamble\} on its first line, so no preamble text lives inside the chapter. The retained file sources/stacks-more-algebra/retained-preamble.tex is a mechanical, line-for-line copy of the project's single shared preamble as staged in the snapshot; nothing was added, removed or redefined and no package was added.\par\medskip
\noindent\textit{Editorial note.} The source URL is the source-level Stacks Project URL, as required when more than one card is produced from one source; the per-claim tag is carried by claim\_key.\par\medskip
\noindent\textit{Editorial note.} Two internal cross-references appear inside the counted proof span, \textbackslash{}ref\{lemma-compare-representatives-perfect\} at line 27910 and \textbackslash{}ref\{lemma-ideals-generated-by-minors\} at line 27960. Both point to lemmas of this same chapter that lie outside the counted span; they are recorded as navigation pointers only, and both applications together with the surrounding computation stay inside the counted proof span.\par\medskip
\noindent\textit{Allowed standard prerequisites (editorial).}\begin{itemize}
\item Ideals I\_i of minors of an augmented map and the behaviour of these ideals for block matrices. invoked through the cited Lemma lemma-ideals-generated-by-minors at proof line 27960 and applied at proof lines 27961-27963. Gives the formula I\_\{r\_i+1\}(T) = I\_\{r\_i+1\}(T\_1) + I\_\{r\_i\}(T\_1) that collapses the two ideals in the case j = i.
\item Localization at primes to test an equality of ideals. used at proof lines 27908-27909. Reduces the general statement to the local case in which the comparison with a trivial complex can be made.
\item The trivial two-term complex concentrated in two adjacent degrees with identity differential. defined at proof lines 27913-27917. Is the summand Q\textasciicircum{}bullet whose addition changes I\_i only by a power of f, which is exactly what the exponents m and n record.
\end{itemize}\medskip
\ifdefined\hypertarget\hypertarget{sourcepointer-1}{}\fi
\noindent\textit{Source reference.} Original lemma 1: sources/stacks-more-algebra/pinned-original-source.tex, source lines 19628--19628. Same-chapter pointer invoked at proof line 27910 to reduce two complexes representing the same derived object to the case where one is a summand of the other; the reduction, the induction the author states is omitted, and every subsequent case distinction stay inside the counted proof span.\par\medskip
\ifdefined\hypertarget\hypertarget{sourcepointer-2}{}\fi
\noindent\textit{Source reference.} Original lemma 2: sources/stacks-more-algebra/pinned-original-source.tex, source lines 1326--1326. Same-chapter pointer supplying the standard formula for the ideals generated by the minors of a block matrix, used at proof line 27960 to collapse I\_\{r\_i+1\}(T) into I\_\{r\_i\}(T\_1); the computation that consumes it, including the values I\_0(T\_2) = A, I\_1(T\_2) = A and I\_p(T\_2) = 0 for p >= 2, stays inside the counted proof span.\par\medskip
\begin{lemma}
\label{lemma-ideal-well-defined}
Let $A$ be a ring and let $f \in A$ be a nonzerodivisor.
Let $M^\bullet$ and $N^\bullet$ be two bounded complexes
of finite free $A$-modules representing the same object of $D(A)$.
Then
$$
f^m I_i(M^\bullet, f) = f^n I_i(N^\bullet, f)
$$
as ideals of $A$ for integers $n, m \geq 0$ such that
$$
m + \sum\nolimits_{j \geq i} (-1)^{j - i}rk(M^j) =
n + \sum\nolimits_{j \geq i} (-1)^{j - i}rk(N^j)
$$
\end{lemma}

\begin{proof}
It suffices to prove the equality after localization at every prime ideal
of $A$. Thus by Lemma \ref{lemma-compare-representatives-perfect}
and an induction argument we omit
we may assume $N^\bullet = M^\bullet \oplus Q^\bullet$
for some trivial complex $Q^\bullet$, i.e.,
$$
Q^\bullet = \ldots \to 0 \to A \xrightarrow{1} A \to 0 \to \ldots
$$
where $A$ is placed in degree $j$ and $j + 1$.
If $j \not = i - 1, i, i + 1$ then we clearly have equality
$I_i(M^\bullet, f) = I_i(N^\bullet, f)$ and $m = n$
and we have the desired equality. If $j = i + 1$ then the maps
$$
(f, d^i) : M^i \to M^i \oplus M^{i + 1}
\quad\text{and}\quad
(f, d^i, 0) : M^i \to M^i \oplus M^{i + 1} \oplus A
$$
have the same nonzero minors hence in this case we also have
$I_i(M^\bullet, f) = I_i(N^\bullet, f)$ and $m = n$. If $j = i$,
then $I_i(M^\bullet, f)$ is the ideal generated by the
$r_i \times r_i$-minors of
$$
(f, d^i) : M^i \to M^i \oplus M^{i + 1}
$$
and $I_i(N^\bullet, f)$ is the ideal generated by
the $(r_i + 1) \times (r_i + 1)$-minors of
$$
(f \oplus f, d^i \oplus 1) :
(M^i \oplus A) \to (M^i \oplus A) \oplus (M^{i + 1} \oplus A)
$$
With suitable choice of coordinates we see that the matrix
of the second map is in block form
$$
T =
\left(
\begin{matrix}
T_1 & 0 \\
0 & T_2
\end{matrix}
\right),
\quad
T_1 = \text{matrix of first map},
\quad
T_2 =
\left(
\begin{matrix}
f \\
1
\end{matrix}
\right)
$$
With notation as in Lemma \ref{lemma-ideals-generated-by-minors}
we have $I_0(T_2) = A$, $I_1(T_2) = A$, $I_p(T_2) = 0$ for $p \geq 2$
and hence $I_{r_i + 1}(T) = I_{r_i + 1}(T_1) + I_{r_i}(T_1) =
I_{r_i}(T_1)$ which means that $I_i(M^\bullet, f) = I_i(N^\bullet, f)$.
We also have $m = n$ so this finishes the case $j = i$.
Finally, say $j = i - 1$. Then we see that $m = n + 1$, thus we have
to show that $fI_i(M^\bullet, f) = I_i(N^\bullet, f)$.
In this case $I_i(M^\bullet, f)$ is the ideal generated by the
$r_i \times r_i$-minors of
$$
(f, d^i) : M^i \to M^i \oplus M^{i + 1}
$$
and $I_i(N^\bullet, f)$ is the ideal generated by
the $(r_i + 1) \times (r_i + 1)$-minors of
$$
(f \oplus f, d^i) : (M^i \oplus A) \to (M^i \oplus A) \oplus M^{i + 1}
$$
With suitable choice of coordinates we see that the matrix
of the second map is in block form
$$
T =
\left(
\begin{matrix}
T_1 & 0 \\
0 & T_2
\end{matrix}
\right),
\quad
T_1 = \text{matrix of first map},
\quad
T_2 =
\left(
\begin{matrix}
f
\end{matrix}
\right)
$$
Arguing as above we find that indeed
$fI_i(M^\bullet, f) = I_i(N^\bullet, f)$.
\end{proof}

\end{document}
